\section{Qué queda de un ciclo con sus ramas}
Consideremos un grafo de factores con un único ciclo simple. Primero no hay
reacción exterior. En las creencias de Bethe de los campos dados, la identidad
de reparametrización expresa el cociente $Z/Z_B$ como la suma de las tablas
de factores multiplicadas y divididas por las marginales repetidas. Para verificarla,
en cada factor $f$, el cociente de la configuración con $i$ respecto de la vacía
es $x_i/q_f$. En el producto de creencias, la ocupación de $i$ aporta así
$x_i^{d_i}/\prod_{f\ni i}q_f$, dividido por $(x_i/(1-x_i))^{d_i-1}$
por los singletons repetidos. El cociente resulta $\zeta_i$ por estación.
En una configuración admitida estos cocientes se multiplican. El peso de la
configuración vacía es $\prod_f q_f/\prod_i(1-x_i)^{d_i-1}=Z_B^{-1}$, por la
misma ecuación de campo y la expresión de $H_B$.

Se suman las hojas hacia el ciclo. En un factor de aridad mayor, desaparecen
las ramas por marginalización y queda su tabla de dos sitios del ciclo. El producto
cíclico de transiciones estocásticas binarias tiene un autovalor uno; el otro es
el producto de los determinantes. Para el arco de $i$ a $k$, el sitio $k$ es cero
si $i$ está ocupado y su probabilidad de ocupación dado $i=0$ es $x_k/(1-x_i)$.
La traza da
\[
 Z/Z_B=1+\theta,\qquad
 \theta=(-1)^d\prod_{i\in C}x_i/(1-x_i),\qquad d=|C|.
\]
Aquí $x_i$ son las creencias de la silla con todas sus ramas. Análogamente la
integral de densidades gaussianas de factores divididas por las repetidas es
$D_{\rm loc}^{-1/2}$. Su precisión es $-H_B''$ pues $C_f^{-1}=-H_{\rm cat}''$.
Se integran primero las ramas. En coordenadas de varianza uno, la regresión
de $k$ sobre $i$ es $\rho_{ik}=-\sqrt{x_ix_k/((1-x_i)(1-x_k))}$.
Cambiar variables a las innovaciones $y_k-\rho_{ik}y_i$ sobre el ciclo introduce
el jacobiano $1-\prod\rho_{ik}=1-\theta$. La integral vale $1/(1-\theta)$.
Se deduce en el origen real y se prolonga por la fórmula holomorfa del apartado 3.
Por consiguiente el error local completo vale $\log(1-\theta^2)$.

La geometría ordinaria impide que dos casillas distintas compartan dos líneas.
El ciclo tiene al menos tres sitios. Al desactivar incidencias sólo puede romperse
o conservar las mismas aristas del ciclo. En cada subcaso su error local es cero
o divisible por $\prod_{i\in C}\zeta_i^2$. Lo mismo vale para la proyección,
que además se anula en cada eje. El puro error local es así
$O_v(j^{-v-d})$ tras proyectar. El cambio por reacción exterior tiene norma
$O_v(s)$ y las mismas compatibilidades de cara. El coste del ciclo con exterior
queda acotado por $O_v(j^{-v-1})$. Cuando hay varios ciclos usaremos únicamente
la holomorfía acotada del apartado 3 y el peso proyectado $O_v(j^{-v})$.
